% Requires in the preamble of thesis.tex:
%   \usepackage{tikz}
%   \usetikzlibrary{positioning, fit, backgrounds, arrows.meta}
\begin{figure}[H]
\centering
\label{fig:code_structure}
\hyphenpenalty=10000\exhyphenpenalty=10000 % no hyphenation in the narrow boxes
\begin{tikzpicture}[
  font=\scriptsize,
  boxc/.style={
    draw=black!55, rounded corners=2pt, fill=white,
    align=left, inner sep=3pt
  },
  boxstep/.style={boxc, text width=6.4cm},
  boxdata/.style={boxstep, fill=green!8, draw=green!45!black},
  boxside/.style={boxc, text width=2.9cm, fill=black!3},
  boxwide/.style={boxside, text width=3.6cm},
  boxout/.style={boxc, text width=3.7cm, fill=blue!4},
  arrflow/.style={-{Stealth[length=2mm]}, draw=black!65, line width=0.6pt},
  arrghost/.style={arrflow, dashed, draw=black!45},
  tag/.style={
    anchor=south west, fill=#1, text=white,
    font=\scriptsize\bfseries, inner sep=2pt, rounded corners=2pt
  },
  filetag/.style={tag=black!55}
]

% ============ 1. user-supplied description of the system ============
\node[boxc, text width=2.9cm] (geom)
  {\texttt{ColumnGeometry}\\ area, height, nozzles};
\node[boxc, text width=2.9cm, right=2mm of geom] (mat)
  {\texttt{BreederMaterial}\\ salt and its properties};
\node[boxc, text width=2.9cm, below=2mm of geom.south west, anchor=north west] (oper)
  {\texttt{OperatingParameters}\\ $T$, $P_\mathrm{top}$, gas flow, breeding};
\node[boxc, text width=2.9cm, right=2mm of oper] (sparg)
  {\texttt{SpargingParameters}\\ closure relations to enforce};

\begin{scope}[on background layer]
  \node[fit=(geom)(mat)(oper)(sparg), inner sep=2.5mm,
        draw=purple!60, fill=purple!5, rounded corners=3pt] (g1) {};
\end{scope}

% ============ 2. resolution of the closure relations ============
\node[boxstep, below=7mm of g1] (resolve)
  {\texttt{SimulationInput.from\_parameters(}\\
   \texttt{\quad geom, mat, oper, sparg)}\\[1mm]
   recursive graph search (\texttt{find\_in\_graph}): every parameter still
   missing is taken from the user input, or else from the default correlation};

% ============ 3. resolved input of the 1D-ARD model ============
\node[boxdata, below=7mm of resolve] (siminput)
  {\texttt{SimulationInput}\\[1mm]
   scalars: \texttt{height}, \texttt{area}, \texttt{temperature}, \texttt{K\_s},
   \texttt{rho\_l}, \texttt{E\_g}, \texttt{E\_l}, \texttt{Q\_T}\\
   profiles $z \mapsto$: \texttt{P\_l}, \texttt{P\_g}, \texttt{eps\_g},
   \texttt{a}, \texttt{u\_g}, \texttt{h\_l}\\
   \texttt{.graph}: intermediate parameters and their dependencies};

% ============ 4. solve ============
\node[boxstep, below=7mm of siminput] (solve)
  {\texttt{Simulation(sim\_input, t\_final,}\\
   \texttt{\quad dispersion\_on, constant\_profiles)}\\
   \texttt{.solve(dt, dx)}\\[1mm]
   1D-ARD weak form assembled and solved with FEniCSx/DOLFINx
   (P1 elements, implicit time stepping, Newton solver)};

% ============ 5. results ============
\node[boxdata, below=7mm of solve] (results)
  {\texttt{SimulationResults}\\[1mm]
   \texttt{times}, \texttt{c\_T2\_profiles}, \texttt{y\_T2\_profiles},
   \texttt{n\_T2\_salt\_series}, and the fields listed in
   \texttt{Simulation.exports}};

% ============ 6. post-processing ============
\node[boxout, below=13mm of results] (post)
  {\texttt{summarize\_decay()}\\ fitted $\tau$ against its analytical prediction};
\node[boxout, left=3mm of post] (save)
  {\texttt{exports\_to\_csv()}\\ \texttt{to\_json(sections)}\\ \texttt{to\_pickle()}};
\node[boxout, right=3mm of post] (anim)
  {\texttt{create\_animation()}\\ interactive $c_{T_2}(z,t)$ and $P_{T_2}(z,t)$};

% ============ side: library of closure relations ============
\node[boxside, left=5mm of resolve] (corr)
  {\texttt{all\_closures}\\[0.5mm]
   \texttt{Correlation} and \texttt{Profile} objects: defaults, or a custom
   one passed in the input};

% ============ side: analytical diagnostics ============
\node[boxside, left=5mm of siminput] (diag)
  {\texttt{get\_Pi\_number()}\\ \texttt{get\_tau\_ave()}\\ \texttt{get\_Bo()}\\
   \texttt{get\_c\_T2\_SS()}\\ \texttt{dx\_from\_Pe()}\\[0.5mm]
   regime check, and sizing of \texttt{dt} and \texttt{dx}};

% ============ side: ready-made inputs ============
\node[boxwide, right=5mm of g1] (examples)
  {\texttt{get\_sim\_input\_LIBRA\_Pi()}\\ \texttt{get\_sim\_input\_standard()}\\[0.5mm]
   pre-built cases, bypassing steps 1 and 2};

% ============ flow ============
\draw[arrflow] (g1) -- (resolve);
\draw[arrflow] (resolve) -- (siminput);
\draw[arrflow] (siminput) -- (solve);
\draw[arrflow] (solve) -- (results);
\draw[arrflow] (corr.east) -- (resolve.west);
\draw[arrflow] (siminput.west) -- (diag.east);
\draw[arrflow] (diag.south) |- ([yshift=2mm]solve.west);
\draw[arrghost] (examples.south) |- (siminput.east);

\draw[arrflow] (results.south) -- ++(0,-4mm) -| (save.north);
\draw[arrflow] (results.south) -- ++(0,-4mm) -| (post.north);
\draw[arrflow] (results.south) -- ++(0,-4mm) -| ([xshift=5mm]anim.north);

% ============ step numbers and source files ============
\node[tag=purple!60]      at (g1.north west)       {1 -- describe the system \textbar{} \texttt{simulation\_input.py}};
\node[filetag]            at (resolve.north west)  {2 -- resolve the closures \textbar{} \texttt{simulation\_input.py}};
\node[tag=green!45!black] at (siminput.north west) {3 -- model input \textbar{} \texttt{simulation\_input.py}};
\node[filetag]            at (solve.north west)    {4 -- solve \textbar{} \texttt{ard\_model.py}};
\node[tag=green!45!black] at (results.north west)  {5 -- results \textbar{} \texttt{ard\_model.py}};
\node[tag=blue!55]        at (save.north west)     {6 \textbar{} \texttt{ard\_model.py}};
\node[tag=blue!55]        at (post.north west)     {\texttt{postprocess.py}};
\node[tag=blue!55]        at (anim.north west)     {\texttt{animation.py}};
\node[filetag]            at (corr.north west)     {\texttt{closure.py}};
\node[filetag]            at (diag.north west)     {\texttt{simulation\_input.py}};
\node[filetag]            at (examples.north west) {\texttt{example\_cases.py}};

\end{tikzpicture}
\caption{High-level structure of the \texttt{sparging} package, and the order in which its API is used to write a simulation}
\end{figure}
